\chapter{Fundamentação Teórica e Trabalhos Relacionados}
\label{ch:fundamentacao}

% Meta editorial: 12--17 páginas. Narrativa: CIL -> interferência -> famílias
% de métodos -> contraste -> destilação contrastiva -> TagFex -> ponto de
% intervenção -> literatura relacionada -> lacuna.

\section{Aprendizado contínuo e incremental de classes}

\subsection{Definição e notação do cenário incremental}
% Definir amostras, rótulos, etapas, classes disjuntas, memória de exemplares,
% espaço de saída acumulado e avaliação após cada etapa.

\subsection{Protocolos e informação de tarefa}
% Diferenciar Task-IL, Domain-IL e Class-IL; task-aware, task-agnostic na
% inferência e task-free/boundary-free no treinamento. Declarar o protocolo
% efetivamente usado pelo TagFex e por este trabalho.

\section{Esquecimento, interferência e estratégias de mitigação}

\subsection{Estabilidade, plasticidade e fontes de erro}
% Relacionar esquecimento, interferência, deriva de representação, viés do
% classificador e desbalanceamento. Não atribuir toda queda a um único mecanismo.

\subsection{Repetição, regularização e expansão arquitetural}
% Revisão focada de rehearsal/herding, destilação, expansão dinâmica e métodos
% híbridos; explicitar custos e pressupostos.

\subsection{Alinhamento e classificação por médias de classe}
% Preparar weight alignment, protótipos e NME usados no pipeline de referência.

\section{Aprendizado contrastivo e destilação contrastiva}

\subsection{Embeddings, similaridade e pares contrastivos}
% Normalização L2, cosseno, temperatura, duas views, âncora, positivo, negativos
% e estrutura da matriz 2B x 2B.

\subsection{InfoNCE e fluxo de gradientes}
% Derivar a perda por âncora e interpretar cada termo. Explicar o que continua
% contribuindo para o gradiente, sem concluir antecipadamente que todo gradiente
% residual seja prejudicial.

\subsection{Destilação contrastiva em aprendizado incremental}
% Distinguir contraste aluno--aluno de destilação aluno--professor e explicar
% por que esse segundo caso é o foco científico do problema.

\subsection{Ponte para os ajustes não essenciais}
% Revisar a motivação da literatura e formular a possibilidade de relações já
% preservadas continuarem ativas. A definição e solução pertencem ao Cap. 3.

\section{TagFex como método de referência}

\subsection{Arquitetura e visão geral do pipeline}
% Expansão de atributos específicos à tarefa, ramo agnóstico, projetor,
% preditor e cabeças. Inserir figura: TagFex -> destilação contrastiva -> ponto
% exato em que a ANT intervém.

\subsection{Primeira etapa e etapas incrementais}
% Explicar expansão, congelamento, professor, duas views, replay e diferenças
% entre o primeiro treinamento e os seguintes.

\subsection{Termos da função objetivo}
% Reunir classificação principal, cabeça/perda auxiliar, contraste,
% destilação contrastiva, classificação de transferência e KL condicional.
% Interpretar coeficientes e o fator automático de destilação.

\subsection{Memória, herding, NME e alinhamento de pesos}
% Explicar quando cada operação ocorre e como participa da avaliação.

\subsection{Uma etapa incremental passo a passo}
% Pseudocódigo/fluxograma alinhado a train() e _compute_incremental_loss() de
% tagfex_clean.py. Encerrar identificando a chamada contrastiva modificada.

\section{Ponto de intervenção e lacuna científica}
% Mostrar em que termo do TagFex as relações contrastivas são construídas, por
% que o critério de ativação é relevante e o que o método-base não controla.
% Não apresentar ainda todas as decisões aGlobal/aLocal/aSymFull.

\section{Trabalhos relacionados}

\subsection{Métodos de CIL diretamente comparáveis}
% TagFex, DER, iCaRL e apenas baselines efetivamente comparáveis.

\subsection{Seleção contrastiva, margens e preservação de representações}
% Revisão curta das linhas mais próximas ao mecanismo proposto.

\subsection{Síntese e transição para a proposta}
% Responder: o que existe, o que cobre, o que falta e como essa lacuna motiva a
% ANT. Terminar apontando diretamente ao Capítulo~\ref{ch:proposta}.
